\section*{Discussion}
A cortical region's response to local optogenetic inhibition was approximately linear and time-invariant on average, but single trials varied with brain state and ongoing activity.
A proportional--integral controller that assumed this linear, time-invariant model tracked the reference and reduced trial-to-trial variability on average by correcting deviations from it.
Its single-trial performance depended on brain state, consistent with disturbances that vary with state.
Replayed on one session's trials, a model-predictive controller that forecast the disturbance lowered the error further.

The impulse experiments set our design choices.
The trial-averaged response grew monotonically with laser power, and a low-order linear model fitted each session's average (Figure~\ref{fig:figure2}B--D), so the controller assumed linearity.
The fitted models transferred only inconsistently between sessions (Figure~\ref{fig:figure2}E), so we recalibrated the feedforward gain in every session.

Single trials also varied with state (Figure~\ref{fig:figure2}G), but what counts as brain state differs between studies.
It can be the activity at stimulus onset, which shaped optogenetic responses in macaque sensorimotor cortex~\citep{bryan2025temporal} and enters our analysis as the initial deviation.
It can also be a variable that changes the dynamics themselves, such as movement and arousal~\citep{Niell2010,Vinck2015,McGinley2015waking} or the synchrony of ongoing activity~\citep{Poulet2008,Harris2011}, which our slow-wave power measures track.
These variables shape the dynamics but are hard to predict, so we measured them post hoc, though identifying them in real time could inform the controller, as we plan next.
Sorted by movement or slow-wave power, single trials departed from their average about as much as stimulus-free activity did from its own, with similar state dependence (Figure~\ref{fig:state_quartile}), consistent with a fixed response riding on ongoing activity~\citep{Arieli1996,Shimaoka2019}.
State that adds an offset is a disturbance that feedback subtracts, whereas state that changes gain is a model error that feedback must tolerate.
Models that treat state as a parameter of the dynamics switch between linear systems~\citep{linderman2017rslds} or model dynamics within a regime, such as the subspace relevant to a behavior~\citep{sani2021modeling}.
Input--output models fitted per state would separate offset from gain, but would require precise identification of each state switch and, for control, ideally its forecast.
A real-time decoder of discrete cortical states, such as a hidden semi-Markov model of LFP spectra in mouse somatosensory cortex~\citep{weiss2024realtime}, labels the state without modelling the response to stimulation, so it could select among interpretable per-state models, although its estimates lag offline ones by about 120--240\,ms.
All such models must contend with inputs they cannot observe, as cortical activity tracks uninstructed movements~\citep{musall2019single,stringer2019spontaneous} and drive from other areas~\citep{zagha2013motor,harris2013topdown} as well as its own dynamics.

We therefore simplified, treating the trial-averaged system as stationary (Assumption~\ref{assm:trial_avg}).
A single linear model is then enough for control if, within a session, the laser's average effect on the readout is well defined and monotonic.
The distribution of states is unknown, and over-sampling one state shifts the average toward that state's mean response but not beyond the range of the per-state means.
That range, unmeasured here, loosely bounds the design model's error.
Averages therefore need enough trials, and in three sessions the cross-trial variance of spontaneous activity changed by only 1--2\% with the number of trials averaged (Figure~\ref{fig:spatial_spread}C).
A nonlinear model is the usual alternative, but linear ones often predict brain activity, including responses to stimulation, about as well~\citep{Nozari2024,sani2024dissociative,Yang2021}, and the richer the model of the dynamics, the more must be known to control with it.
Both kinds forecast ongoing activity only briefly, and stimulus-blind, univariate forecasts of spontaneous mouse wide-field activity stay informative for about 1.5\,s, with deep models modestly ahead of a linear autoregressive one~\citep{lu2025benchmarking}.
Authors of a zebrafish whole-brain benchmark expect an irreducible error floor~\citep{lueckmann2025zapbench}.
Feedback partly opposes such unforecastable inputs through the error they cause.
A model for feedback design need only capture how the laser moves the readout~\citep{gevers2005identification}, and a low-order linear model makes stability, gain limits and bandwidth computable~\citep{astrom2008feedback}.

Proportional--integral control needs only the sign of the laser's effect, which was consistently suppressive (Figure~\ref{fig:figure2}B; Figure~\ref{fig:figure3}B).
The simplest optimal regulator of a system must model it~\citep{conant1970every}, and ours modelled only its average.
The feedforward gain approximately inverted the average response, and the integrator, an internal model of a constant signal rather than of cortex, drives a constant offset toward zero within its clamp~\citep{francis1976internal}.
Monotonicity gives each attainable reference one steady-state light drive, which low-gain integral control of an exponentially stable linear plant reaches despite a saturating actuator and input delay~\citep{logemann2000absolute}.
Session-to-session variability therefore limits the gains, not whether some correctly signed gain stabilizes the loop~\citep{morari1985robust}.
With this minimal model, feedback lowered error in every session and, on average, trial-to-trial variability (Figure~\ref{fig:figure3}H,I).
Against the contralaterally predicted disturbance~\citep{Shimaoka2019,Ye2026}, feedback lowered residual energy in all 11 retained sessions (two excluded for large closed-loop offsets), mostly by shrinking the standing offset left by open-loop calibration, while fluctuations fell less (Figure~\ref{fig:figure4}G; Table~\ref{tab:stat_fig4g}).

Earlier closed-loop optogenetic controllers held neuronal firing rates in culture and in rodent thalamus~\citep{newman2015optogenetic,bolus2018design,bolus2021state}, clamped single-neuron calcium signals~\citep{Zhang2018} or phase-locked light to oscillations~\citep{nicholson2018analogue,kanta2019closed,Zaaimi2022}.
In anesthetized rat thalamus, proportional--integral feedback held firing rates although its fitted model predicted open-loop responses poorly, lowered Fano factors below open loop in five cells, and could not lower rates above target with an excitatory opsin~\citep{bolus2018design}, the mirror image of our one-sided limit.
Changing brain state was anticipated as a disturbance~\citep{bolus2021state} or a change in the plant~\citep{grosenick2015closed}, whereas here we related residual closed-loop error to measured state.

Poorly regulated trials may reflect activity outside the controller's model, so per-trial error probes brain state: slow 1--4\,Hz power explained part of the 1--3\,s error, and movement none (Figure~\ref{fig:figure4}B,C).
Open-loop error rose with movement while closed-loop error stayed flat.
Slow, large-amplitude 1--4\,Hz fluctuations were reduced, but closed-loop error still rose with their power as much as open-loop error (Figure~\ref{fig:figure4}D; Table~\ref{tab:fig4_state_stats}).
An effective delay of 37--61\,ms (loop delay at most 47\,ms and about half that on average, Figure~\ref{fig:figure1}F, plus the 14\,ms half-sample hold) caps the feedback bandwidth at about 2.6--4\,Hz by the crossover criterion $\omega_c\tau \lesssim 1$\,rad~\citep{skogestad2005multivariable}.
Cortical lag (Figure~\ref{fig:figure5}H) probably lowers the unmeasured crossover further.
Large slow excursions may also saturate the command, which can suppress activity but not raise it.
A closed- to open-loop error spectrum and time at the ceiling would separate delay from saturation.
The band-power window spans stimulation, which the loop shapes, and motion and slow power are anticorrelated~\citep{McGinley2015waking}, so state effects are associations.
Where state changes the plant itself, the alternative to a single time-invariant model is a context-dependent model whose controller switches, or schedules its gains, with state~\citep{rugh2000research,grosenick2015closed}.
Adaptive deep brain stimulation switches on when a measured biomarker crosses a threshold~\citep{little2013adaptive}, rather than adjusting stimulation continuously, and optogenetic feedback in mouse thalamus paired a first-order model with an adaptive estimator that compensated for mismatch online~\citep{bolus2021state}.

Model-predictive control is a modest step beyond proportional--integral control, adding prediction to a linear model that now includes the rest of cortex and treating the one-sided laser as a constraint~\citep{mayne2000constrained}.
Preview, which predicts only the reference, already cancelled the cortical phase lag that feedback left (Figure~\ref{fig:figure5}H).
The contralateral hemisphere predicted much of the site's spontaneous activity (Figure~\ref{fig:figure4}F).
Yet in one session, 100 cortical sites added little to a 200\,ms forecast of the site beyond its own past and the laser (Figure~\ref{fig:mpc_direct}H; Figure~\ref{fig:supp_varx}E), so drive from the rest of cortex was observable but largely unforecastable at short leads.
A linear model without subcortical signals may understate what is forecastable.
In a replay of one session, the controller lowered the error that feedback left by about 40\%, and by about 80\% given the future residuals (Figure~\ref{fig:mpc_direct}C--F).
What remained was mostly activity its model could not predict.
The replayed cortex shares the controller's model class and responds to the laser identically in every state, so these values bound, rather than measure, live performance, as for model-based designs for electrical stimulation tested in simulation~\citep{Yang2018,Yang2021}.
A live controller should keep integral action~\citep{pannocchia2003disturbance}.

The data comprise four impulse sessions in three mice, 15 separate controller sessions in four mice (11 in the disturbance analysis) and three sinusoidal sessions in one mouse.
In every impulse session the average response grew monotonically with power and was predicted by a linear model (Figure~\ref{fig:figure2}B--D), so with a comparable preparation it should be predictable enough for control.
All sessions inhibited one somatosensory site in awake, head-fixed mice, leaving excitation, other areas and task performance untested, though the preparation reaches all dorsal cortex~\citep{Matveev2024}.
The loop fed back on fluorescence without hemodynamic correction, which cannot run within one control interval, so closed-loop measures describe the regulated signal rather than purely neural activity~\citep{valley2020separation}.
Inhibition spreads toward the adjacent readout (Figure~\ref{fig:spatial_spread}A,B;~\citealp{li2019spatiotemporal}), so the loop holds a readout kernel, not a whole area, and partly regulates that spread.
Interleaved trials limit, but cannot remove, drift in brain state between conditions.
Gains were tuned empirically, and the low-gain guarantee assumes a fixed plant, so stability rests on the observation that no session showed sustained oscillation.
There was no sham, opsin-free light-only, random-timing or proportional-only condition.
Opsin-independent light effects would therefore enter the identified plant, and proportional action also shrinks an offset, so the role of integral action in removing it is inferred, not isolated.

We present a procedure for precise feedback control that makes cortical activity follow a chosen reference without an explicit model of its dynamics.
Characterize the input--output relationship with brief stimuli at several intensities and check that the average response is monotonic.
If a low-order linear model fits that average and the trial average is stable in the number of trials (Assumption~\ref{assm:trial_avg}), proportional--integral control with feedforward compensation, tuned to the experiment, can be used.
Nothing in this procedure is specific to wide-field imaging or optogenetics, as it requires only a graded stimulus that moves the readout monotonically and stable trial averages.
It should therefore carry over to two-photon imaging, electrophysiology, electrocorticography or functional MRI, and to electrical stimulation, within each readout's bandwidth.
Model-predictive control may add precision where the disturbance is partly predictable.
A site that tracks a chosen reference also lets its influence on dorsal cortex, striatum~\citep{Peters2021} or behavior~\citep{ZatkaHaas2021} be tested at defined activity levels.
As in the voltage clamp and the optoclamp~\citep{newman2015optogenetic}, the laser command needed to track the reference should then report the drive that the loop opposes, a prediction we did not test.
The linear design is also an investigative tool, as the error it leaves shows which brain states its model misses and so guides context-dependent models and switching controllers.
